% Study guide for learning target 1A. Build: python3 units/ecology/study-guides/build.py 1A/guide
% Macros: latex/preamble.tex (\drawfig: figures/ in this folder, \bookcrop) and latex/guide.sty (\guidetitle, tbpage, \tbnote).
\documentclass[11pt,letterpaper]{article}
\newcommand{\sgroot}{../}
% Shared preamble for the study guides and mock exams (built by ../build.py).
% \sgroot is the path from the document's folder to study-guides/ (empty for guides, ../ for exams).
\providecommand{\sgroot}{}
\usepackage[margin=0.8in,headheight=15pt]{geometry}
\usepackage{fontspec}
\setmainfont{texgyrepagella}[Extension=.otf,UprightFont=*-regular,BoldFont=*-bold,ItalicFont=*-italic,BoldItalicFont=*-bolditalic]
\setsansfont{texgyreheros}[Extension=.otf,UprightFont=*-regular,BoldFont=*-bold,ItalicFont=*-italic,BoldItalicFont=*-bolditalic]
\usepackage{amsmath,amssymb,booktabs,longtable,array,calc,enumitem,xcolor,graphicx}
\usepackage{tikz}
\usepackage{pgfplots}\pgfplotsset{compat=1.18}
\usetikzlibrary{calc,arrows.meta,positioning}
\usepackage[most]{tcolorbox}
\usepackage{fancyhdr,titlesec,needspace,etoolbox,environ}
\usepackage{hyperref}

\definecolor{forest}{HTML}{24594B}
\definecolor{pale}{HTML}{EFF5F1}
\definecolor{keyc}{HTML}{C98A0B}   % key idea: amber
\definecolor{vocabc}{HTML}{2F6FA8} % vocabulary: blue
\definecolor{tipc}{HTML}{3F8A3A}   % make it make sense: green
\definecolor{watchc}{HTML}{B8452F} % watch out: red
\definecolor{linkc}{HTML}{7B4C9A}  % connection: purple
\hypersetup{colorlinks=true,linkcolor=forest,urlcolor=vocabc}

\pagestyle{fancy}\fancyhf{}
\renewcommand{\headrulewidth}{0pt}
\fancyfoot[C]{\small\sffamily\thepage}
\setlength{\parindent}{0pt}\setlength{\parskip}{5pt}
\setlength{\emergencystretch}{2em}
\setlist[itemize]{leftmargin=1.3em,itemsep=2pt,topsep=2pt}
\setlist[enumerate]{leftmargin=1.6em,itemsep=2pt,topsep=2pt}
\renewcommand{\arraystretch}{1.2}
\providecommand{\tightlist}{\setlength{\itemsep}{0pt}\setlength{\parskip}{0pt}}
\providecommand{\pandocbounded}[1]{#1}
\def\LTcaptype{none}

\titleformat{\section}{\Large\sffamily\bfseries\color{forest}}{}{0pt}{}[{\color{forest}\titlerule[0.8pt]}]
\titleformat{\subsection}{\large\sffamily\bfseries\color{forest}}{}{0pt}{}
\titleformat{\subsubsection}{\normalsize\sffamily\bfseries}{}{0pt}{}
\titlespacing*{\section}{0pt}{14pt}{8pt}
\titlespacing*{\subsection}{0pt}{12pt}{4pt}
\setcounter{secnumdepth}{0}

% Block quotes (textbook excerpts) as shaded boxes.
\renewenvironment{quote}{\begin{tcolorbox}[breakable,enhanced,colback=pale,colframe=forest,boxrule=0pt,leftrule=2.5pt,arc=0pt,
  left=8pt,right=8pt,top=4pt,bottom=4pt,before skip=4pt,after skip=6pt]\small}{\end{tcolorbox}}

% ------------------------------------------------------------------ annotated textbook pages
% Landscape letter page: the scanned page on the left (full height), notes on the right.
\newlength{\portraitW}\setlength{\portraitW}{8.5in}
\newlength{\portraitH}\setlength{\portraitH}{11in}
\def\pgmargin{21.6}  % pt, top/left margin around the scan
\def\pgheight{568.8} % pt, printed height of the scan (8.5in - 2 margins)
\newcommand{\setpagesize}[2]{%
  \global\pdfpagewidth=#1\global\pdfpageheight=#2%
  \global\paperwidth=#1\global\paperheight=#2}

\newcommand{\circlednum}[2]{\tikz[baseline=(c.base)]\node[circle,fill=#1,text=white,font=\scriptsize\sffamily\bfseries,inner sep=1.2pt,minimum size=11pt](c){#2};}

% \textbookpage{image}{page width pt}{page height pt}{lesson}{book page}{intro}{highlights}{notes}
% Notes (\addnote) are typeset into boxes first, then laid out: each sits level with its highlight
% when there is room, is pushed down to avoid the note above, and is pulled back up from the page bottom.
\newcount\nnotes
\newbox\hdrbox
\newdimen\prevbottom \newdimen\limitdim \newdimen\tmpdim
\def\notegap{7pt}
\newcommand{\addnote}[5]{% color, number, label, anchor (scan pt, <0 = none), text
  \global\advance\nnotes by 1
  \expandafter\ifx\csname nb@\the\nnotes\endcsname\relax\expandafter\newbox\csname nb@\the\nnotes\endcsname\fi
  \global\expandafter\setbox\csname nb@\the\nnotes\endcsname=\hbox{\tikz\node[notebox=#1]{\notehead{#1}{#2}{#3}#5};}%
  \expandafter\xdef\csname na@\the\nnotes\endcsname{#4}%
  \expandafter\xdef\csname nc@\the\nnotes\endcsname{#1}}
\newcommand{\nb}[1]{\csname nb@#1\endcsname}
\newcommand{\layoutnotes}{%
  \global\prevbottom=\dimexpr\pgmargin pt+\ht\hdrbox+\dp\hdrbox\relax
  \ifnum\nnotes>0
    \foreach \k in {1,...,\the\nnotes}{%
      \edef\a{\csname na@\k\endcsname}%
      \tmpdim=\dimexpr\prevbottom+\notegap\relax
      \ifdim \a pt<0pt\else
        \pgfmathsetlength{\dimen0}{\pgmargin+\a*\s-7}%
        \ifdim\dimen0>\tmpdim \tmpdim=\dimen0 \fi
      \fi
      \expandafter\xdef\csname ny@\k\endcsname{\the\tmpdim}%
      \global\prevbottom=\dimexpr\tmpdim+\ht\nb{\k}+\dp\nb{\k}\relax}%
    \global\limitdim=\dimexpr612pt-32pt\relax
    \foreach \k in {\the\nnotes,...,1}{%
      \tmpdim=\csname ny@\k\endcsname
      \ifdim\dimexpr\tmpdim+\ht\nb{\k}+\dp\nb{\k}\relax>\limitdim
        \tmpdim=\dimexpr\limitdim-\ht\nb{\k}-\dp\nb{\k}\relax
        \expandafter\xdef\csname ny@\k\endcsname{\the\tmpdim}%
      \fi
      \global\limitdim=\dimexpr\tmpdim-\notegap\relax}%
    \ifdim\limitdim<\dimexpr\pgmargin pt+\ht\hdrbox+\dp\hdrbox-\notegap\relax
      \typeout{WARNING: notes overflow on textbook page \currentbookpage}\fi
  \fi}
\newcommand{\drawnotes}{%
  \ifnum\nnotes>0
    \foreach \k in {1,...,\the\nnotes}{%
      \edef\a{\csname na@\k\endcsname}\edef\c{\csname nc@\k\endcsname}\edef\y{\csname ny@\k\endcsname}%
      \node[anchor=north west,inner sep=0pt] at ($(O)+(\notex pt,-\y)$) {\usebox{\nb{\k}}};
      \ifdim \a pt<0pt\else
        \draw[\c,line width=0.6pt] ($(scan.north east)+(0,{-\a*\s pt})$) -- ++(6pt,0) -- ($(O)+(\notex pt-1pt,-\y-7pt)$);
      \fi}%
  \fi}
\newcommand{\textbookpage}[8]{%
  \clearpage\setpagesize{11in}{8.5in}\thispagestyle{empty}%
  \def\currentbookpage{#5}%
  \pgfmathsetmacro{\s}{\pgheight/#3}%
  \pgfmathsetmacro{\pgright}{\pgmargin+#2*\s}%
  \pgfmathsetmacro{\notex}{\pgright+20}%
  \pgfmathsetmacro{\notew}{792-\notex-\pgmargin}%
  \pgfmathsetmacro{\notetw}{\notew-10}%
  \setbox\hdrbox=\vbox{\hsize=\notew pt\parskip=0pt
    {\raggedright\sffamily\bfseries\color{forest}#4\par}\vspace{1pt}
    {\sffamily\footnotesize\color{black!55}Miller \& Levine Biology, textbook page #5\par}\vspace{-1pt}
    {\color{forest}\rule{\notew pt}{0.6pt}}\par\vspace{2pt}
    \ifx\relax#6\relax\else{\small\raggedright #6\par}\fi}%
  \global\nnotes=0 #8\layoutnotes
  \null
  \begin{tikzpicture}[remember picture,overlay]
    \coordinate (O) at (current page.north west);
    \node[anchor=north west,inner sep=0pt,draw=black!25] (scan) at ($(O)+(\pgmargin pt,-\pgmargin pt)$)
      {\includegraphics[height=\pgheight pt]{#1}};
    \node[anchor=north west,inner sep=0pt] at ($(O)+(\notex pt,-\pgmargin pt)$) {\usebox{\hdrbox}};
    \node[anchor=south east,font=\scriptsize\sffamily,text=black!50,inner sep=0pt] at ($(O)+(792pt-\pgmargin pt,-603pt)$)
      {Part 1 · the textbook, annotated \quad·\quad Miller \& Levine Biology (2019), p.\,#5 \quad·\quad page \thepage};
    #7
    \drawnotes
  \end{tikzpicture}%
  \clearpage\setpagesize{\portraitW}{\portraitH}}

% Highlight one text line on the scan: \hl{color}{x0}{y0}{x1}{y1} in scan points (y down).
\newcommand{\hl}[5]{%
  \fill[#1,opacity=0.22] ($(scan.north west)+({(#2-1.5)*\s pt},{-(#3-1)*\s pt})$) rectangle
                         ($(scan.north west)+({(#4+1.5)*\s pt},{-(#5+1)*\s pt})$);}
% Numbered badge just left of a highlight.
\newcommand{\badge}[4]{%
  \node[circle,fill=#1,text=white,font=\scriptsize\sffamily\bfseries,inner sep=1pt,minimum size=10pt]
    at ($(scan.north west)+({#3*\s pt-7pt},{-#4*\s pt})$) {#2};}
\newcommand{\notehead}[3]{{\sffamily\bfseries\footnotesize\circlednum{#1}{#2}\ \color{#1}\MakeUppercase{#3}}\\[1pt]}
\tikzset{notebox/.style={text width=\notetw pt,inner sep=4pt,fill=#1!7,draw=#1!55,line width=0.4pt,
  font=\small,align=left,rounded corners=1.5pt}}

% ------------------------------------------------------------------ figures in the guide text
% Caption used under both book figures and new diagrams.
\newcommand{\figcaption}[2]{\par\smallskip{\small\sffamily\raggedright #1\par}%
  \ifx\relax#2\relax\else{\scriptsize\sffamily\color{black!55}Source: #2\par}\fi}
% \bookfig{png}{width fraction}{caption}{source}: a figure cropped from one of the books.
\newcommand{\bookfig}[4]{\par\Needspace{8\baselineskip}\medskip
  \begin{minipage}{\linewidth}\centering
    \fcolorbox{black!15}{white}{\includegraphics[width=#2\linewidth]{#1}}
    \figcaption{#3}{#4}
  \end{minipage}\par\medskip}
% Several book figures side by side: \bookfigrow{\bookfigcol{w}{png}{height}{caption}{source}\hfill ...}
\newcommand{\bookfigcol}[5]{\begin{minipage}[t]{#1\linewidth}\centering
    \fcolorbox{black!15}{white}{\includegraphics[height=#3,width=0.95\linewidth,keepaspectratio]{#2}}
    \figcaption{#4}{#5}\end{minipage}}
\newcommand{\bookfigrow}[1]{\par\Needspace{10\baselineskip}\medskip\noindent#1\par\medskip}
% New diagrams drawn for the guide: \begin{guidefig}...\end{guidefig}{caption}
\newenvironment{guidefig}[1]{\par\Needspace{10\baselineskip}\medskip\begin{minipage}{\linewidth}\centering\def\guidefigcap{#1}}
  {\figcaption{\guidefigcap}{}\end{minipage}\par\medskip}
\tikzset{
  gbox/.style={draw=#1,fill=#1!8,rounded corners=3pt,align=center,font=\small\sffamily,inner sep=5pt},
  garr/.style={-{Stealth[length=2.4mm]},line width=1pt,draw=#1},
  glab/.style={font=\scriptsize\sffamily,fill=white,inner sep=1.5pt,align=center},
}

% ------------------------------------------------------------------ figures shared by guides and exams
% \drawfig{name}: a TikZ/pgfplots figure in figures/name.tex
\newcommand{\drawfig}[1]{\input{\sgroot figures/#1}}
% \bookcrop[width]{id}{book}{page}{x0 y0 x1 y1}{label}{caption}: a figure cropped from a book page.
% book: C = Campbell Biology, ML = Miller & Levine Biology; page = printed page number;
% crop box in points from the page's top-left corner. build.py makes build/figs/<id>.png from these arguments.
\newcommand{\bookname}[1]{\ifstrequal{#1}{C}{Campbell Biology}{Miller \& Levine Biology}}
\newcommand{\bookcrop}[7][0.7]{\bookfig{\sgroot build/figs/#2.png}{#1}{#7}{\bookname{#3}, #6, p.\,#4}}

\usepackage{../latex/guide}
\hypersetup{pdftitle={1A Study Guide},pdfauthor={Prepared for Shawn}}
\InputIfFileExists{../build/anchors/1A.tex}{}{}

\begin{document}
\guidetitle{1A}{1A Study Guide: Academic and Personal Skills}
  {The development of academic and personal skills occurs through deliberate practice.}

\section{How to use this guide}

You don't need the textbook to use this guide. Every textbook page that matters for 1A is copied into \textbf{Part 1}, in order, with the important sentences highlighted. Each highlight has a numbered badge that matches a note in the column beside the page. The notes explain hard words, answer the textbook's questions, and connect the page to what your teacher asks about in class.

\begin{enumerate}
  \item \textbf{Read Part 1 first.} Read each textbook page, then its notes. Don't skip the notes on the reading checks and review questions, because they give you model answers.
  \item \textbf{Then read Parts 2–5} (after the textbook pages): the other sources the textbook leaves out, a one-page summary, the vocabulary explained in detail, and full answers to the teacher's study questions.
  \item \textbf{Test yourself.} Cover the answers in Part 5 and try each study question out loud before you check.
\end{enumerate}

Notes are color-coded:
\notelegend

\section{What this guide covers}

\textbf{Where to read}

\begin{longtable}{@{}>{\raggedright\arraybackslash}p{0.24\linewidth}>{\raggedright\arraybackslash}p{0.26\linewidth}>{\raggedright\arraybackslash}p{0.07\linewidth}>{\raggedright\arraybackslash}p{0.30\linewidth}@{}}
\toprule
Source & Section & Pages & What it covers \\
\midrule
\endhead
Class textbook: \emph{Miller \& Levine Biology} (2019) & Lesson 1.1 What Is Science? & 8–14 & observation, inference, hypothesis, data (quantitative and qualitative) \\
Class textbook: \emph{Miller \& Levine Biology} (2019) & Lesson 1.2 Science in Context & 15–18 & scientific attitudes, collaboration, feedback, peer review \\
Reference: \emph{Campbell Biology} & Ch. 1, Concept 1.3 & 16–17 & observation, data, inductive and deductive reasoning \\
Teacher's slides & Biology Class Basics (KJ\_PH) & slides 3–8 & class policies, formative and summative assessment, academic skills \\
Teacher's slides & Deductive vs Inductive Reasoning Overview & slides 3–6 & induction and deduction, with practice examples \\
Teacher's worksheets & Observation and Inference; Geotastic Inference Game & — & Sherlock Holmes activity, types of data, a prediction-testing game \\
\bottomrule
\end{longtable}

The unit guide lists no textbook lesson for 1A (1B onward list one). The lessons above were chosen because they cover 1A's vocabulary and study questions. The textbook does not cover class policies or induction and deduction. Those answers use the teacher's slides and Campbell.

\clearpage

% ================================================================ Part 1: the textbook, annotated
% Each tbpage is one landscape page: the textbook page on the left, the notes on the right.

\begin{tbpage}{8}{Lesson 1.1 What Is Science?}
  {This lesson answers ``what is science?'' Most of 1A comes from here: observation, inference, and data. The list of words at the top left is the lesson's vocabulary.}
  \tbnote{vocab}{Lesson vocabulary}{observation inference hypothesis}{theory}
    {For 1A, focus on \textbf{observation}, \textbf{inference}, and \textbf{data}. Hypothesis, controlled experiment, variables, and control group come back in 1B.}
  \tbnote{key}{}{scientific knowledge is always changing}{}
    {Scientific knowledge changes when new evidence appears. That is why science is a skill you practice, not a list of facts to memorize.}
  \tbnote{watch}{}{Scientists either understand}{reject that explanation.}
    {Scientists don't ``believe'' in ideas. They \textbf{accept or reject} explanations based on evidence. Avoid writing ``scientists believe\dots{}'' in answers.}
  \tbnote{vocab}{science}{The term science is usually defined}{through this process.}
    {Science has two meanings: a \textbf{process} (the way we find things out) and the \textbf{body of knowledge} that process produces.}
  \tbnote{key}{}{Most importantly, science is a process}{reject those explanations.}
    {The process in four steps: \textbf{observe and ask questions $\rightarrow$ make a testable explanation $\rightarrow$ gather and analyze data $\rightarrow$ support or reject the explanation.} Worth memorizing.}
\end{tbpage}

\begin{tbpage}{9}{Lesson 1.1 What Is Science?}
  {How science is different from other ways of explaining the world, and what science is trying to do.}
  \tbnote{key}{}{First, science deals only with the natural world.}{not belief.}
    {Three features of science: (1) it studies only the \textbf{natural} world; (2) it collects information in an \textbf{orderly} way and looks for patterns and cause and effect; (3) its explanations come from \textbf{evidence}, not belief.}
  \tbnote{tip}{}{universal natural laws}{}
    {The same natural laws apply to something huge (a hurricane) and something tiny (a cell). That's why one experiment can teach us about many situations.}
  \tbnote{key}{Goals of science}{One goal of science is to provide}{about natural events.}
    {Science aims to (1) give \textbf{natural, testable explanations}, (2) understand \textbf{patterns} in nature, and (3) make useful \textbf{predictions}.}
  \tbnote{link}{}{How do chimpanzees interact}{}
    {Each question can be answered with data. ``Ideal temperature for basil'' gives \textbf{quantitative} data (numbers). ``How chimps interact'' gives mostly \textbf{qualitative} data (descriptions of behavior).}
\end{tbpage}

\begin{tbpage}{10}{Lesson 1.1 What Is Science?}
  {Science is uncertain and always improving. The ``scientific method'' is really a flexible way of thinking that you already use.}
  \tbnote{key}{}{science rarely}{exciting!}
    {Science doesn't \textbf{prove} things. It finds the best explanation the evidence supports. This connects to study question 4: science relies on induction, which is never 100\% certain.}
  \tbnote{tip}{}{And please don't believe them}{those ideas.}
    {The book tells you to be \textbf{skeptical} of the book! Ask \emph{how} scientists know something, not just \emph{what} they know.}
  \tbnote{tip}{Everyday science}{Suppose your family}{call a mechanic!}
    {Each guess (``dead battery?'', ``out of gas?'') is a \textbf{hypothesis}, and each check (turn the key, look at the gauge) is a \textbf{test}. When a test fails, you reject that idea and try the next one.}
  \tbnote{watch}{}{There isn't a single, cut-and-dried}{}
    {There is \textbf{no single fixed ``scientific method''} with the same steps every time. Scientists follow a general style called \emph{scientific methodology}.}
  \tbnote{key}{}{Scientific methodology involves}{drawing conclusions.}
    {Five parts: \textbf{observe and ask questions $\rightarrow$ form hypotheses $\rightarrow$ do controlled experiments $\rightarrow$ collect and analyze data $\rightarrow$ draw conclusions.} Figure 1-2 on the next page walks through all five.}
  \tbnote{tip}{Reading check: answer}{Construct an Explanation How is scientific}{}
    {Scientific knowledge is based on \textbf{evidence} from observation and experiment, it must be \textbf{testable}, and it is always open to \textbf{revision} when new evidence appears. Other kinds of knowledge (opinion, belief, tradition) don't have to meet those rules.}
\end{tbpage}

\begin{tbpage}{11}{Lesson 1.1 What Is Science?}
  {Figure 1-2 follows one real study from start to finish. Match each numbered step to the 1A vocabulary.}
  \tbnote{vocab}{Step 1: observation}{OBSERVING AND ASKING QUESTIONS}{}
    {The \textbf{observation} is that grass is taller in some places. It is something you can see and describe. The question it leads to: \emph{why?}}
  \tbnote{vocab}{Step 2: inference}{INFERRING AND HYPOTHESIZING}{}
    {``Something limits growth'' is an \textbf{inference}, an interpretation of the observation using background knowledge. Narrowing it to one testable cause (nitrogen) makes it a \textbf{hypothesis}.}
  \tbnote{link}{}{DESIGNING CONTROLLED EXPERIMENTS}{}
    {The plots are alike in everything except nitrogen. Only one thing differs between groups: that's a \textbf{controlled experiment} (1B vocabulary).}
  \tbnote{vocab}{Step 4: data}{COLLECTING DATA}{}
    {Growth rates, plant sizes, and leaf chemistry are all \textbf{quantitative} data (measured numbers). A note like ``the grass was lying flat'' would be qualitative.}
  \tbnote{watch}{}{ANALYZING CONCLUSIONS}{}
    {Look at the wording: the hypothesis was \textbf{supported}, not \emph{proved}. The graph shows the +N plots (red) growing taller than the control (blue) over time.}
\end{tbpage}

\begin{tbpage}{12}{Lesson 1.1 What Is Science?}
  {The most important page for 1A. It defines observation, inference, and the two kinds of data.}
  \tbnote{vocab}{observation}{Scientific investigations begin with observation}{careful, orderly way.}
    {An \textbf{observation} is noticing and \textbf{describing} carefully, using your senses or instruments. It says \emph{what is there}, not what it means.}
  \tbnote{tip}{}{The right kind of observation leads}{answered) before.}
    {Good observations lead to good \textbf{questions}. Sherlock Holmes is famous because he notices details other people miss.}
  \tbnote{vocab}{inference}{An inference is a logical interpretation}{already know.}
    {An \textbf{inference} goes \emph{beyond} the observation to explain it, using what you already know. It can be wrong even when the observation is right.}
  \tbnote{key}{}{Inference, combined with a creative imagination}{by experimentation.}
    {The chain: \textbf{observation $\rightarrow$ inference $\rightarrow$ hypothesis $\rightarrow$ test.} A hypothesis is an inference stated so it can be tested.}
  \tbnote{vocab}{quantitative data}{Quantitative data are numbers}{counting or measuring.}
    {\textbf{Numbers} from counting or measuring. *Quanti*tative = *quanti*ty. Examples: 47 plants, 30 cm tall, 22\,°C.}
  \tbnote{vocab}{qualitative data}{Qualitative data are descriptive}{cannot usually be measured.}
    {\textbf{Descriptions} instead of numbers. *Quali*tative = *quali*ty (what kind). Examples: upright or sideways grass, yellow leaves, a plot disturbed by clam diggers.}
\end{tbpage}

\begin{tbpage}{13}{Lesson 1.1 What Is Science?}
  {How scientists make their data trustworthy, and what to do when an experiment isn't possible.}
  \tbnote{tip}{}{Statistical analysis helps}{from controls.}
    {\textbf{Statistics} check whether a difference is real or just chance. If the nitrogen plots were only 1 cm taller, that might be luck.}
  \tbnote{key}{}{The larger the sample size}{experimentals and controls.}
    {\textbf{More data make conclusions more reliable.} That's the idea behind the Geotastic game: more advances (more observations) should give better location guesses (better inferences).}
  \tbnote{key}{}{Data analysis may lead to conclusions that support or refute}{experimental design.}
    {Conclusions \textbf{support or refute} (reject) a hypothesis. Often the result is ``partly right,'' which leads to a revised hypothesis and a better experiment. Science improves the same way you do: test, get feedback, revise.}
  \tbnote{tip}{}{Not all hypotheses can be tested by experiments.}{}
    {Some questions (like animal behavior in the wild) are answered by careful \textbf{field observations} instead of experiments. Jane Goodall's chimpanzee research is a famous example.}
  \tbnote{tip}{Reading check: answer}{Describe What is the difference between}{}
    {\textbf{Quantitative} data are numbers from counting or measuring (height = 30 cm). \textbf{Qualitative} data are descriptions of characteristics that usually can't be measured (the grass is leaning sideways).}
  \tbnote{tip}{Analyzing Data: answers}{Analyze Graphs According to the circle graph}{}
    {1. \textbf{Fruits} (49\% of the diet; leaves are 38\%). 2. With less fruit, siamangs would have less of their main food. They might eat more leaves, lose weight, or have fewer young. A hypothesis: \emph{if fruit decreases, then siamang populations will decline.} The graph is \textbf{quantitative} data.}
\end{tbpage}

\begin{tbpage}{14}{Lesson 1.1 What Is Science?}
  {What a scientific theory is, plus the lesson review. The notes give model answers to every review question.}
  \tbnote{watch}{}{When most people say}{}
    {In everyday speech, ``theory'' means a guess. In science it means the \textbf{opposite}: a well-tested, highly reliable explanation.}
  \tbnote{vocab}{theory}{In science, the word theory applies to}{accurate predictions.}
    {A \textbf{theory} is a tested, highly reliable explanation that ties together many observations and well-supported hypotheses, and that makes accurate predictions. Examples: cell theory, evolution.}
  \tbnote{key}{}{But remember that no theory is absolute truth.}{}
    {Even a strong theory can be \textbf{revised} if new evidence appears. Again: science supports, it doesn't prove.}
  \tbnote{tip}{Review 1–3: answers}{In your own words, define the term science.}{}
    {\textbf{1.} Science is an organized way of using evidence to explain and predict events in the natural world, plus the knowledge that process builds. \textbf{2.} A hypothesis gives the experiment its purpose: it says what to change (independent variable) and what result to expect. \textbf{3.} A hypothesis is a single testable explanation. A theory is a well-tested explanation that unifies many observations and hypotheses.}
  \tbnote{tip}{Review 4–5: answers}{Form a Hypothesis You observe mold}{}
    {\textbf{4.} Example: \emph{the moldy side got more moisture} (or less light, or touched a mold source). \textbf{5.} Hypothesis: warmer water increases goldfish activity. Independent variable: water temperature (e.g., 15, 20, 25\,°C). Dependent variable: gill beats per minute. Control group: fish at normal tank temperature. Keep fish size, tank size, and food the same, and use several fish per group.}
\end{tbpage}

\begin{tbpage}{15}{Lesson 1.2 Science in Context}
  {Lesson 1.2 shows science as a messy loop, not a straight line, and as something done by a community.}
  \tbnote{key}{}{The testing of ideas is the heart}{process of science.}
    {Scientists question, observe, look for evidence, \textbf{share ideas}, and analyze data. All of these are part of science, not just experiments.}
  \tbnote{key}{}{Notice that the parts of the process do not}{proceed.}
    {Real science \textbf{loops back} on itself. Results send you back to new questions. This is the same loop as deliberate practice in 1A: try, get feedback, fix, try again.}
  \tbnote{vocab}{bias}{bias}{}
    {A \textbf{bias} is ``a personal, rather than scientific, point of view for, or against, something'' (p. 20). Controlled experiments and peer review help keep bias out of results.}
\end{tbpage}

\begin{tbpage}{16}{Lesson 1.2 Science in Context}
  {Where scientific ideas come from, and the habits of mind that good scientists (and good students) use.}
  \tbnote{link}{}{Ideas in science can arise}{}
    {Scientists start by observing, asking questions, \textbf{talking with colleagues}, and reading earlier work. That is \textbf{collaboration} (a 1A academic word).}
  \tbnote{key}{}{Curiosity, skepticism, open-mindedness, and creativity}{new problems.}
    {The four \textbf{scientific attitudes}. Know all four, and be ready to give an example of each.}
  \tbnote{vocab}{skepticism}{Skepticism Scientists and engineers should be skeptics}{without evidence.}
    {Being \textbf{skeptical} means questioning ideas and asking for evidence before accepting them. It's not the same as refusing to believe anything.}
  \tbnote{vocab}{open-mindedness}{Open-Mindedness Scientists and engineers should be}{}
    {Being \textbf{open-minded} means being willing to accept evidence and ideas that disagree with what you expected. Pairs well with skepticism: question everything, but change your mind when the evidence says so.}
\end{tbpage}

\begin{tbpage}{17}{Lesson 1.2 Science in Context}
  {How scientists work together and check each other's work.}
  \tbnote{vocab}{collaboration}{Scientists often collaborate in groups}{scientific community.}
    {Scientists \textbf{collaborate}: they work in teams and share with other teams. Results only count once they are shared with the scientific community.}
  \tbnote{vocab}{feedback}{Communication is an important part of science.}{}
    {Figure 1-6: scientists give each other \textbf{feedback} through peer review, discussion, and replication. Your formative assessments do the same job for you.}
  \tbnote{tip}{}{Scientists also share their work with the general public}{}
    {Scientists also share informally, through news articles, TV, and blogs, especially when results could help society.}
  \tbnote{tip}{Reading check: answer}{Describe Describe a}{}
    {Example: ``Cracking your knuckles causes arthritis.'' A skeptic asks what evidence supports it. (Studies comparing knuckle-crackers with non-crackers have not found more arthritis.)}
\end{tbpage}

\begin{tbpage}{18}{Lesson 1.2 Science in Context}
  {Peer review: how the scientific community checks new research before accepting it.}
  \tbnote{vocab}{peer review}{Peer Review Scientists share findings}{called peer review.}
    {\textbf{Peer review:} other experts in the same field check a paper for mistakes, bias, and fraud before it is published. It's feedback and assessment for scientists.}
  \tbnote{watch}{}{Peer review does not guarantee}{scientific community.}
    {Peer review doesn't \textbf{guarantee} a paper is right. It shows the work meets the community's standards.}
  \tbnote{key}{}{These logical and important}{}
    {Every answer raises \textbf{new questions} (does nitrogen limit mangroves too?). Science keeps going.}
  \tbnote{tip}{Reading check: answer}{Why is it necessary for other scientists to evaluate}{}
    {Other experts can catch mistakes, bias, or fraud the authors missed, and can \textbf{replicate} the work to see if they get the same result. That makes the findings more trustworthy.}
  \tbnote{link}{}{Replicating}{}
    {Quick Lab: if your directions are unclear, others can't \textbf{replicate} your work. This is the 1A academic skill of writing clearly and concisely.}
\end{tbpage}

% ================================================================ Part 1 (continued)
\section{Part 1 (continued). Other sources for this target}

The class textbook pages are reproduced in full above. These excerpts cover the rest: the reference book \emph{Campbell Biology} and anything the class textbook leaves out.

\subsection{Campbell Biology, Chapter 1, Concept 1.3}

The teacher's \emph{Observation and Inference} worksheet asks questions about this passage (``The scientific process begins with?'', ``Recorded observations are called?'').

\textbf{Observation} (p. 16)

\begin{quote}
Biology, like other sciences, begins with careful observations. In gathering information, biologists often use tools such as microscopes, precision thermometers, or high-speed cameras that extend their senses or facilitate careful measurement. \dots{} In exploring nature, biologists also rely heavily on the scientific literature, the published contributions of fellow scientists.
\end{quote}

\textbf{Data} (p. 16)

\begin{quote}
\textbf{Recorded observations are called data.} Put another way, data are items of information on which scientific inquiry is based. The term \emph{data} implies numbers to many people. But some data are qualitative, often in the form of recorded descriptions rather than numerical measurements. For example, Jane Goodall spent decades recording her observations of chimpanzee behavior \dots{} Quantitative data are generally expressed as numerical measurements and often organized into tables and graphs.
\end{quote}

\textbf{Inductive reasoning} (p. 17)

\begin{quote}
Collecting and analyzing observations can lead to important conclusions based on a type of logic called inductive reasoning. \textbf{Through induction, we derive generalizations from a large number of specific observations.} ``The sun always rises in the east'' is one example. Another biological example is the generalization ``All organisms are made of cells,'' which was based on two centuries of microscopic observations \dots{}
\end{quote}

\textbf{Deductive reasoning} (p. 17)

\begin{quote}
While induction entails reasoning from a set of specific observations to reach a general conclusion, \textbf{deductive reasoning involves logic that flows in the opposite direction, from the general to the specific.} \dots{} In the scientific process, deductions usually take the form of a prediction of results that will be found if a particular hypothesis (premise) is correct. \dots{} This deductive testing takes the form of ``If . . . then'' logic. In the case of the desk lamp example: If the burnt-out bulb hypothesis is correct, then the lamp should work if you replace the bulb with a new one.
\end{quote}

\textbf{A hypothesis is never proved} (p. 17)

\begin{quote}
Although a hypothesis can never be proved beyond all doubt, testing it in various ways can significantly increase our confidence in its validity.
\end{quote}

% ================================================================ Part 2. Summary
\section{Part 2. Summary}

\textbf{The big idea of 1A.} Science is a \emph{skill}, not a list of facts, and skills improve through deliberate practice. Deliberate practice means: try it, get feedback, fix the gap, try again. That is the same loop scientists use, where they observe, explain, test, get feedback from peers, and revise. This first target teaches two sets of skills together:

\drawfig{practice-loop}

\begin{enumerate}
  \item \textbf{How science thinks.}
  \begin{itemize}
    \item Science begins with \textbf{observation}: noticing and describing carefully, using senses and instruments.
    \item Recorded observations are \textbf{data}. \textbf{Quantitative} data are numbers (counts, measurements). \textbf{Qualitative} data are descriptions (color, behavior, shape).
    \item From observations, scientists make \textbf{inferences}, logical interpretations based on what they already know. An inference plus creativity can become a testable \textbf{hypothesis}.
    \item \textbf{Induction} moves from many specific observations to a general conclusion (``every swan I've seen is white $\rightarrow$ all swans are white''). Its conclusion is \emph{probably} true, never guaranteed.
    \item \textbf{Deduction} moves from a general statement to a specific prediction (``if nitrogen limits growth, then adding nitrogen should make grass taller''). If the premises are true, the conclusion must be true.
    \item Science uses both. Induction builds the general idea, and deduction turns it into an if–then prediction you can test. Because the general idea came from induction, science never fully ``proves'' anything. Evidence makes an explanation more or less likely.
    \item Science is a community effort: \textbf{collaboration}, \textbf{feedback}, and \textbf{peer review} catch mistakes and improve ideas.
  \end{itemize}
\end{enumerate}

\begin{enumerate}
  \item \textbf{How this class works.}
  \begin{itemize}
    \item The \textbf{KUDOS} and \textbf{Unit Guide} tell you what to learn: vocabulary, where to find information, and study questions.
    \item \textbf{Formative} assessments (homework, practice tests, CFAs) are practice. They don't affect your grade and exist to give you feedback.
    \item \textbf{Summative} assessments (quizzes, unit exams, final) measure what you learned and count toward your grade.
    \item Academic skills you are expected to practice: note taking, reading comprehension, \textbf{annotation}, clear and concise writing, academic vocabulary, CER argumentation, and study and test-taking strategies.
  \end{itemize}
\end{enumerate}

\textbf{One-sentence version:} observe carefully, label what you \emph{saw} separately from what you \emph{think it means}, reason from evidence, and use feedback to keep improving, both in science and in this class.

\bookcrop[0.78]{c-process}{C}{18}{30 262 612 762}{Figure 1.23}
  {The process of science is a loop, not a straight line. Testing ideas (center) connects to exploring (top), community feedback (right), and benefits to society (left). Deliberate practice works the same way.}

% ================================================================ Part 3. Biology vocabulary
\section{Part 3. Biology vocabulary}

\subsection{annotation}
\begin{itemize}
  \item \textbf{Definition:} Notes you write directly on a text, diagram, or data table while reading. You mark key ideas, define unfamiliar words, ask questions, and make connections.
  \item \textbf{Source:} Biology Class Basics, slide 8 (listed as an academic skill). The textbook doesn't define it, though its Reading Tool boxes ask you to annotate.
  \item \textbf{What good annotation looks like:}
  \begin{itemize}
    \item Underline or highlight the \emph{one} key sentence in a paragraph, not the whole paragraph. (In Miller \& Levine, the key sentence often starts with a small key icon.)
    \item Write a short note in the margin in your own words: ``inference = interpretation, not observation.''
    \item Put ``?'' next to anything confusing and write the question down.
    \item Draw arrows to connect ideas: ``$\rightarrow$ same as the marsh grass example.''
  \end{itemize}
  \item \textbf{Why it matters:} It turns reading into active thinking, which is deliberate practice. You also remember more, and your notes make review faster before a test.
  \item \textbf{Common mistake:} highlighting everything. If everything is marked, nothing stands out.
\end{itemize}

\drawfig{annotation-example}

\subsection{formative (formative assessment)}
\begin{itemize}
  \item \textbf{Definition:} An assessment used \emph{during} learning to check understanding and give feedback. It is practice, not a final score.
  \item \textbf{In this class:} homework, practice tests, and CFAs (common formative assessments). According to the Class Basics slides, formative assessments \textbf{do not affect your grade} and are a ``chance for feedback and improvement.''
  \item \textbf{Memory hook:} \emph{form*ative helps *form} your learning, like clay still being shaped.
  \item \textbf{How to use it:} treat each formative as information. Every mistake tells you exactly what to practice before the real test.
\end{itemize}

\subsection{summative (summative assessment)}
\begin{itemize}
  \item \textbf{Definition:} An assessment at the \emph{end} of a period of learning that measures what you learned and counts toward your grade.
  \item \textbf{In this class:} quizzes, unit exams, and the final exam. Quizzes/exams are 45\% of the grade and the final is 20\%.
  \item \textbf{Memory hook:} \emph{sum*mative is the *sum} of what you learned.
  \item \textbf{Formative vs. summative:} the same kind of question can appear on both. What differs is the \emph{purpose}: to practice and get feedback (formative) or to measure and grade (summative).
\end{itemize}

\drawfig{formative-summative}

\subsection{quantitative / qualitative (data)}
\begin{itemize}
  \item \textbf{Quantitative data:} ``numbers obtained by counting or measuring'' (M\&L p. 12). Examples: number of plants per plot, blade length in cm, temperature in\,°C, a heart rate of 72 beats per minute.
  \item \textbf{Qualitative data:} ``descriptive and involve characteristics that cannot usually be measured'' (M\&L p. 12). Examples: grass growing upright or sideways, leaf color is yellow, the chimp groomed another chimp (Jane Goodall's notes, Campbell p. 16).
  \item \textbf{Memory hook:} *quanti*tative = *quanti*ty (how many, how much); *quali*tative = *quali*ty (what kind, what it's like).
  \item \textbf{Same object, both kinds:} a microscope can count blood cells (quantitative) or show the steps of cell division (qualitative). The worksheet uses this example.
  \item \textbf{Watch out:} a description with a number-like word is still qualitative if nothing was counted or measured. ``Many plants'' is qualitative; ``47 plants'' is quantitative.
\end{itemize}

\bookcrop[0.5]{c-goodall}{C}{16}{312 586 588 762}{Figure 1.21}
  {Qualitative data: Jane Goodall recorded chimpanzee behavior as written descriptions and sketches.}

\subsection{data}
\begin{itemize}
  \item \textbf{Definition:} ``Recorded observations are called data'' (Campbell p. 16). Data are the information scientific inquiry is based on, collected with senses or instruments.
  \item \textbf{Two types:} quantitative and qualitative (above).
  \item \textbf{Role in science:} data \emph{support or reject} a hypothesis (M\&L p. 8). Scientists often organize data in tables and graphs and use statistics to check whether a result is real or just random chance (M\&L p. 13; Campbell p. 16).
  \item \textbf{Grammar note:} \emph{data} is plural (the singular is \emph{datum}), so scientists write ``the data show \dots{}''
  \item \textbf{Good data are:} recorded carefully, repeated (replicated), and taken from a large enough sample. The larger the sample size, the more reliable the analysis (M\&L p. 13).
\end{itemize}

\bookcrop[0.55]{c-mice-data}{C}{20}{30 48 406 602}{Figure 1.25}
  {Quantitative data: the percentage of mouse models attacked by predators, shown as a bar graph. Camouflaged models were attacked much less.}

\subsection{inference}
\begin{itemize}
  \item \textbf{Definition:} ``a logical interpretation based on what scientists already know'' (M\&L p. 12). An inference goes \emph{beyond} what was directly observed to explain it.
  \item \textbf{Observation vs. inference:}
\end{itemize}

  | Observation (what you detect) | Inference (what you conclude it means) |   |---|---|   | The grass in location A is 90 cm tall; in location B it is 30 cm. | Something in location B is limiting growth. |   | The road signs are in Portuguese. (Geotastic) | This is probably Brazil or Portugal. |   | There is a tan line on the man's wrist, but not his hand. (Sherlock) | He has recently been somewhere sunny. |   | The lamp won't light. | The bulb is probably burnt out. |

\begin{itemize}
  \item \textbf{Key point:} an inference can be wrong even when the observation is right. That's why scientists turn inferences into hypotheses and test them.
  \item \textbf{Test tip:} if a statement contains ``because,'' ``probably,'' ``must be,'' or a cause, it's usually an inference.
\end{itemize}

\bookcrop[0.9]{c-mice}{C}{19}{62 556 614 766}{Figure 1.24}
  {Observation: beach mice are light and inland mice are dark, each matching its surroundings. Inference: the coloring camouflages them from predators. That inference became a hypothesis, tested in Figure 1.25.}

% ================================================================ Part 4. Academic vocabulary
\section{Part 4. Academic vocabulary}

These are general academic words, not just biology words. You'll see them on tests and in instructions.

\subsection{assess / assessment}
\begin{itemize}
  \item \textbf{Meaning:} to \emph{assess} is to judge or measure the quality, amount, or level of something. An \emph{assessment} is the tool or event that does the judging: a test, quiz, checkup, or lab report.
  \item \textbf{In class:} formative and summative are the two kinds of assessment.
  \item \textbf{Example sentence:} ``The teacher used a checkup to \emph{assess} whether students completed the homework.''
  \item \textbf{In science:} scientists assess evidence to decide whether it supports a hypothesis. Peer reviewers assess each other's papers.
\end{itemize}

\subsection{collaboration}
\begin{itemize}
  \item \textbf{Meaning:} working together with others toward a shared goal. (Verb: \emph{collaborate}; adjective: \emph{collaborative}.)
  \item \textbf{In science:} ``Scientists often collaborate in groups and communicate with other research groups'' (M\&L p. 17). Discussing with colleagues, sharing data, replication, and peer review are all forms of collaboration (M\&L Figure 1-6).
  \item \textbf{In class:} lab partners, group discussion, sharing class data in a Google Sheet (Geotastic).
  \item \textbf{Caution:} collaboration means building understanding together, not copying. You still need to be able to explain the idea on your own.
\end{itemize}

\bookcrop[0.42]{c-community}{C}{18}{366 545 612 762}{Figure 1.23 (detail)}
  {Community analysis and feedback: scientists collaborate through peer review, replication, publication, and discussion.}

\subsection{commitment}
\begin{itemize}
  \item \textbf{Meaning:} a promise or dedication to keep doing something, especially when it's hard. (Verb: \emph{commit}.)
  \item \textbf{In class:} 15–30 minutes of \emph{focused} homework per night (Class Basics slide 7), coming to tutorial when stuck, and doing formative practice even though it isn't graded.
  \item \textbf{In science:} Jane Goodall's decades of observing chimpanzees (Campbell p. 16), and Darwin spending many years assembling his hypotheses into a theory (M\&L p. 14).
  \item \textbf{Link to 1A:} deliberate practice only works with commitment, meaning repeated, focused effort over time.
\end{itemize}

\subsection{effective}
\begin{itemize}
  \item \textbf{Meaning:} successful in producing the intended result. (Noun: \emph{effectiveness}.)
  \item \textbf{Don't confuse:} \emph{effective} (it works) and \emph{efficient} (it works without wasting time or energy). Rereading notes for an hour can be \emph{inefficient} and \emph{ineffective}. Quizzing yourself for 20 minutes is usually more effective.
  \item \textbf{Example sentence:} ``Annotating the reading was an \emph{effective} way to prepare for the quiz.''
  \item \textbf{In science:} a controlled experiment is an effective design because changing only one variable lets you tell which variable caused the result (M\&L p. 12).
\end{itemize}

\subsection{feedback}
\begin{itemize}
  \item \textbf{Meaning:} information about your performance that tells you what is working and what to improve.
  \item \textbf{In class:} formative assessments are a ``chance for feedback and improvement'' (Class Basics slide 6). Comments on assignments and retake opportunities are also feedback.
  \item \textbf{In science:} ``Community Analysis and Feedback'' (M\&L p. 17). Peer reviewers give feedback that catches ``mistakes, oversights, unfair influences, or fraud'' (M\&L p. 18).
  \item \textbf{How to use it:} (1) read the feedback, (2) find the specific gap, (3) practice that exact skill, (4) check again. That loop \emph{is} deliberate practice.
  \item \textbf{Heads-up for 1B:} in biology, \emph{feedback} also means something different: a \emph{feedback loop} in the body (negative and positive feedback) that keeps conditions stable. Same word, related idea: output information is used to adjust the next step.
\end{itemize}

\subsection{induction / deduction}
\begin{itemize}
  \item \textbf{Induction (inductive reasoning):} reasoning from \textbf{specific observations $\rightarrow$ a general conclusion}. Campbell p. 17: ``we derive generalizations from a large number of specific observations.''
  \begin{itemize}
    \item Example: ``Every time I drop an apple it falls. $\rightarrow$ All dropped apples will fall.'' (teacher's slide 5)
    \item Example: ``Two centuries of microscope observations found cells in every organism $\rightarrow$ All organisms are made of cells.'' (Campbell p. 17)
    \item The conclusion is only \emph{probably} true. New evidence could overturn it.
  \end{itemize}
  \item \textbf{Deduction (deductive reasoning):} reasoning from \textbf{general premises $\rightarrow$ a specific conclusion}. Campbell p. 17: ``from the general to the specific.''
  \begin{itemize}
    \item Example: ``All antibiotics kill bacteria. Penicillin is an antibiotic. Therefore penicillin kills bacteria.'' (teacher's slide 5)
    \item If the premises are true and the logic is valid, the conclusion \textbf{must} be true. It can only be false if a premise is false.
    \item In science, deduction takes the form of an \textbf{if \dots{} then} prediction.
  \end{itemize}
  \item \textbf{Memory hook:} \textbf{In}duction = going \textbf{in}to a big idea from small pieces (specific $\rightarrow$ general). \textbf{De}duction = going \textbf{down} from the big idea to one case (general $\rightarrow$ specific).
\end{itemize}

\drawfig{induction-deduction}

\begin{itemize}
  \item \textbf{Quick check} (teacher's slide 5):
  \begin{itemize}
    \item ``All dolphins are mammals. All mammals have kidneys. Dolphins have kidneys.'' $\rightarrow$ \emph{deductive}
    \item ``Chimps and humans share 99\% of their genetic material. Chimps and humans evolved from a common ancestor.'' $\rightarrow$ \emph{inductive}
    \item ``5, 8, 11, 14, 17, \_\_. The next number is 20.'' $\rightarrow$ \emph{inductive} (a pattern is assumed to continue)
  \end{itemize}
\end{itemize}

\subsection{observation}
\begin{itemize}
  \item \textbf{Meaning:} ``the act of noticing and describing events or processes in a careful, orderly way'' (M\&L p. 12). Also the thing noticed (``my observation was \dots{}'').
  \item \textbf{How we observe:} with the five senses, extended by instruments such as microscopes, thermometers, high-speed cameras, rulers, scales, and pH meters (Campbell p. 16). Scientists also rely on others' published observations, the scientific literature (Campbell p. 16).
  \item \textbf{Good observations are:} specific, factual, and free of interpretation. ``The leaf is yellow with brown edges'' is an observation. ``The leaf is dying'' is an inference.
  \item \textbf{Why it matters:} observation starts every investigation, and ``the right kind of observation leads to asking questions that no one has asked (or answered) before'' (M\&L p. 12).
\end{itemize}

% ================================================================ Part 5. Study questions with answers
\section{Part 5. Study questions with answers}

Try answering each question yourself first, then check.

\subsection{Q1. What are the policies for Biology class?}

\emph{Source: Biology Class Basics slides 3–8. This isn't in the textbook.}

\textbf{How you know what to do}
\begin{itemize}
  \item \textbf{KUDOS} (Know, Understand, Do): a one-page overview of the main topics in the unit.
  \item \textbf{Unit Guide:} more detailed information about what you need to know. It has the important vocabulary, places to find information, and study questions (like this one).
  \item \textbf{Class material:} homework, classwork, lecture notes, and review material. \emph{All} assignments prepare you for the unit exams.
  \item \textbf{Ask a teacher or TA:} if you're not sure, send an email or come to tutorial/office hours.
\end{itemize}

\textbf{Assessments and grades}

\begin{longtable}{@{}>{\raggedright\arraybackslash}p{0.14\linewidth}>{\raggedright\arraybackslash}p{0.22\linewidth}>{\raggedright\arraybackslash}p{0.14\linewidth}>{\raggedright\arraybackslash}p{0.37\linewidth}@{}}
\toprule
Type & Examples & Counts toward grade? & Rules \\
\midrule
\endhead
Formative & practice tests, CFAs, homework & \textbf{No} & Practice. A chance for feedback and improvement. \\
Checkups & cover assigned homework/classwork & Yes, \textbf{15\%} & Not retakeable or made up unless you have an excused absence. Your \textbf{lowest 2 per semester are dropped}. \\
Quality assignments & — & Yes, \textbf{20\%} & — \\
Quizzes and exams & mid-unit or end of unit & Yes, \textbf{45\%} & Quizzes are placeholders. The end-of-unit exam can replace your quiz score. You can \textbf{retake an exam if you score below 80\%}, but you must show homework completion and complete a retake ticket. \\
Final exam & end of course & Yes, \textbf{20\%} & — \\
\bottomrule
\end{longtable}

\drawfig{grade-weights}

\textbf{Homework}
\begin{itemize}
  \item Expect about \textbf{15–30 minutes per night} of \emph{focused} homework (not multitasking).
  \item Homework isn't checked or graded directly, but the \textbf{checkups} assess it. Skipping homework shows up there.
  \item If you've spent \textbf{more than one hour} of focused time and are still stuck, email your teacher and come to the next tutorial.
\end{itemize}

\textbf{Academic skills you'll practice all year:} note taking, reading comprehension, annotation, writing that is clear and concise, correct academic and scientific vocabulary, CER-style scientific argumentation, and study and test-taking strategies.

\subsection{Q2. How will students know what they need to learn and how they can demonstrate their understanding in Biology?}

\textbf{Knowing what to learn.} The teacher tells you directly, in layers:
\begin{enumerate}
  \item The \textbf{KUDOS} gives the one-page big picture of the unit: what you should \emph{know}, \emph{understand}, and be able to \emph{do}.
  \item The \textbf{Unit Guide} breaks the unit into learning targets (1A, 1B, 1C \dots{}). For each one it lists the biology and academic vocabulary, the related classwork (where to find the information), the textbook lessons, and the study questions. If you can answer every study question and use every vocabulary word correctly, you know what the target requires.
  \item \textbf{Class materials} (notes, slides, labs, homework) teach the content. Every assignment prepares you for the exam.
\end{enumerate}

\textbf{Demonstrating understanding.} You show what you know in two ways:
\begin{enumerate}
  \item \textbf{Formative assessments} (homework, practice tests, CFAs): you practice and get \textbf{feedback}. They don't count, so their job is to \emph{reveal gaps} while there's still time to fix them.
  \item \textbf{Summative assessments} (checkups, quizzes, exams, the final): these measure and grade your understanding.
\end{enumerate}

You demonstrate understanding using the academic skills: accurate vocabulary, clear and concise writing, and especially \textbf{CER} (Claim, Evidence, Reasoning). A strong answer makes a claim, supports it with data or observations, and explains \emph{why} the evidence supports the claim.

\textbf{Connecting it to 1A:} the system is built for deliberate practice. Use the unit guide to know the target, practice with formative work, use the feedback to fix gaps, then show mastery on the summative. If you score below 80\% on an exam, the retake policy lets you run the loop again.

\drawfig{learn-demonstrate}

\subsection{Q3. How are observation and inference related to science?}

\textbf{Short answer:} Science begins with observation. Inferences are the explanations scientists build from those observations. A good inference becomes a testable hypothesis, and new observations then test it. Science is a cycle of observing and inferring.

\textbf{In detail:}

\begin{enumerate}
  \item \textbf{Observation comes first.} ``Scientific investigations begin with observation, the act of noticing and describing events or processes in a careful, orderly way'' (M\&L p. 12). Scientists observe with their senses and with instruments that extend their senses, like microscopes and thermometers (Campbell p. 16). Recorded observations are \textbf{data} (Campbell p. 16), either quantitative or qualitative.
\end{enumerate}

\begin{enumerate}
  \item \textbf{Inference interprets the observation.} ``An inference is a logical interpretation based on what scientists already know'' (M\&L p. 12). An observation tells you \emph{what}. An inference proposes \emph{why} or \emph{what it means}. Inferences use prior knowledge, so the same observation can lead different people to different inferences.
\end{enumerate}

\begin{enumerate}
  \item \textbf{Inference leads to a hypothesis.} ``Inference, combined with a creative imagination, can lead to a hypothesis,'' a tentative explanation that can be tested (M\&L p. 12).
\end{enumerate}

\begin{enumerate}
  \item \textbf{More observations test it.} The hypothesis predicts what \emph{should} be observed if it's true. Scientists collect new data through experiments or field observations. The data support or reject the inference.
\end{enumerate}

\drawfig{obs-inf-chain}

\textbf{Worked example: the salt marsh (M\&L Figure 1-2, p. 11)}

\begin{longtable}{@{}>{\raggedright\arraybackslash}p{0.27\linewidth}>{\raggedright\arraybackslash}p{0.60\linewidth}@{}}
\toprule
Step & What happened \\
\midrule
\endhead
Observation & Marsh grass grows taller in some places than others. \\
Inference & Something is limiting growth in the short-grass places. \\
Hypothesis & Marsh grass growth is limited by available nitrogen. \\
Test (new observations) & Add nitrogen to experimental plots, none to control plots, and measure growth. \\
Result & Nitrogen plots grew taller $\rightarrow$ hypothesis supported. \\
\bottomrule
\end{longtable}

\textbf{Worked example: Sherlock Holmes (class worksheet).} Sherlock notices small details, such as a tan line, a worn sleeve, or mud on a shoe. Those are \textbf{observations}, and most of them are \textbf{qualitative} (descriptions, not measurements). From them he draws conclusions about where someone has been or what their job is. Those are \textbf{inferences}. He is effective because his observations are unusually careful and his inferences use a lot of background knowledge. But an inference can be wrong even if the observation is right. That's why science, unlike a detective story, \emph{tests} inferences.

\textbf{Worked example: Geotastic (class activity).} In the game, the clues you see (language on signs, which side of the road cars drive on, vegetation) are observations. Your guess of the location is an inference. The class prediction was that more ``advances'' (more observations) would bring guesses closer to the real location, meaning a smaller distance in km. That is the scientific idea that \textbf{more and better observations lead to more accurate inferences.}

\textbf{Why the distinction matters in science:}
\begin{itemize}
  \item Mixing up observation and inference leads to bad conclusions. You might treat a guess as a fact.
  \item When you write a CER, your \textbf{evidence} must be observations/data, and your \textbf{claim} is an inference. Keeping them separate is what makes the argument convincing.
\end{itemize}

\subsection{Q4. What is the difference between induction and deduction? How does this relate to science?}

\textbf{The difference}

\begin{longtable}{@{}>{\raggedright\arraybackslash}p{0.13\linewidth}>{\raggedright\arraybackslash}p{0.37\linewidth}>{\raggedright\arraybackslash}p{0.37\linewidth}@{}}
\toprule
 & Induction (inductive reasoning) & Deduction (deductive reasoning) \\
\midrule
\endhead
Direction & Specific $\rightarrow$ general & General $\rightarrow$ specific \\
Starts with & Evidence, data, observations & General rules or premises \\
Ends with & A general rule or pattern that explains the evidence & A specific conclusion or prediction \\
Certainty & \textbf{Not guaranteed.} Only \emph{probably} true, based on past experience. & \textbf{Guaranteed} \emph{if} the premises are true and the logic is valid. \\
Teacher's example & ``The sun has always risen. Therefore, the sun will rise tomorrow.'' & ``All dogs have four legs. My pet is a dog. Therefore, my pet has four legs.'' \\
Campbell's example & ``All organisms are made of cells'' (from two centuries of observations) & ``If the burnt-out bulb hypothesis is correct, then the lamp should work if you replace the bulb.'' \\
\bottomrule
\end{longtable}

\begin{itemize}
  \item \textbf{Induction} (Campbell p. 17): ``we derive generalizations from a large number of specific observations.'' It assumes the past predicts the future, and that isn't always true. So an inductive conclusion is only ``probabilistically true based on past experience'' (teacher's slide 4).
  \item \textbf{Deduction} (Campbell p. 17): ``logic that flows in the opposite direction, from the general to the specific.'' The conclusion can only be false if one of the premises is false (teacher's slide 3).
\end{itemize}

\textbf{How it relates to science: science uses both, in a cycle.}

\begin{enumerate}
  \item \textbf{Induction builds hypotheses and theories.} Scientists observe many specific cases and look for a pattern. The hypothesis is ``a rational accounting for a set of observations, based on the available data and guided by inductive reasoning'' (Campbell p. 17).
  \item \textbf{Deduction makes testable predictions.} From the hypothesis, scientists reason: \emph{if} the hypothesis is true, \emph{then} this specific result should happen. ``This deductive testing takes the form of 'If . . . then' logic'' (Campbell p. 17). Example: \emph{If} nitrogen limits marsh grass growth, \emph{then} grass given extra nitrogen should grow taller than grass without it.
  \item \textbf{Experiments check the prediction.} If the result matches, the hypothesis is supported. If not, it's rejected or revised. Then the cycle repeats.
\end{enumerate}

\bookcrop[0.55]{c-desklamp}{C}{17}{345 106 602 346}{Figure 1.22}
  {Deduction in action. Each hypothesis leads to an if \dots{} then prediction (``if the bulb is burnt out, replacing it will fix the lamp''), which a test supports or rejects.}

\textbf{Why science never ``proves'' anything.} The teacher's slide 6 says science ``is not a deductive process, it is an inductive one.'' The general ideas science rests on come from induction, so they can never be 100\% certain:
\begin{itemize}
  \item More supporting data \emph{increases the probability} that a theory is true, but never makes it certain.
  \item New data that goes against the pattern can reduce that probability.
  \item Someone might find a pattern or theory that explains the existing evidence better.
\end{itemize}

The textbooks say the same thing: ``science rarely 'proves' anything'' (M\&L p. 10), and ``a hypothesis can never be proved beyond all doubt'' (Campbell p. 17). Confidence grows as a hypothesis survives more tests, and this is how scientific consensus forms.

\textbf{A connection to the Sherlock worksheet (Q16 on that sheet).} Sherlock is called the ``Master of Deduction,'' but most of what he does is \textbf{induction}. He starts from specific observations (a tan line, a limp) and reaches a general conclusion (``you've been in Afghanistan''). His conclusions are very probable, not guaranteed. So a strong answer is: \emph{No, not really. He mostly reasons from specific observations to a general conclusion, which is inductive reasoning. His conclusions could be wrong if another explanation fits the same clues.} (Arguing the other way is possible when he applies a general rule to a specific case, for example ``people who have been in the tropics are tanned; this man is tanned \dots{}''. Either way, name the direction of reasoning and use a specific example from the clip.)

% ================================================================ Sources
\section{Sources}

\begin{itemize}
  \item Miller, K. R., \& Levine, J. S. \emph{Miller \& Levine Biology}. Pearson, 2019. Lesson 1.1 (pp. 8–14), Lesson 1.2 (pp. 15–18).
  \item Urry, L. A., et al. \emph{Campbell Biology}. Ch. 1, Concept 1.3 (pp. 16–17).
  \item Teacher materials: \emph{Ecology Unit Guide (KJ)}, \emph{Biology Class Basics (KJ\_PH)}, \emph{Deductive vs Inductive Reasoning Overview}, \emph{Observation and Inference}, \emph{Geotastic Inference Game}.
\end{itemize}

\end{document}
